\newpage
\section{Grundlagen} \label{sec:grundlagen}

Dieses Kapitel legt die Begriffe fest, die die folgenden Kapitel voraussetzen. Es ordnet sie in
drei Gruppen: die Schichten, auf denen fremder Code voneinander getrennt wird, die Bausteine, aus
denen die Ausführung auf Kubernetes zusammengesetzt ist, und die beiden Eigenschaften, die die
Forschungsfrage verlangt. Begriffe, die nur an einer Stelle gebraucht werden, erklärt die
jeweilige Stelle.

\subsection{Schichten der Isolation}
\label{sec:grundlagen:isolation}

Ein \emph{Container} grenzt einen Prozess über Mechanismen des Betriebssystemkerns ab, teilt sich
diesen Kern aber mit allen anderen Containern desselben Hosts
\autocite[S.~18]{souppaya_application_2017}. Was daraus für die Angriffsfläche folgt, behandelt
\autoref{sec:angreifermodell}.
In Kubernetes ist der \emph{Pod} die kleinste Einheit, die sich ausbringen lässt. Er fasst einen
oder mehrere Container zusammen, die sich Netz und Speicher teilen und gemeinsam auf einen Knoten
geplant werden \autocite{the_kubernetes_authors_pods_2026}.

\emph{gVisor} schiebt zwischen Container und Host-Kernel eine weitere Schicht. Es ist ein in Go
geschriebener Anwendungskern, der eine Linux-ähnliche Schnittstelle nachbildet. Sein Kernstück,
der Sentry, fängt die Systemaufrufe der Anwendung ab und beantwortet sie selbst, statt sie an den
Host-Kernel durchzureichen. Die Laufzeit \texttt{runsc} setzt die Spezifikation der Open Container
Initiative um und lässt sich dadurch in Kubernetes als Runtime-Klasse eintragen
\autocite{the_gvisor_authors_what_2026}. Weil der Sentry im User Space läuft und nur eine verkleinerte
Oberfläche des Host-Kernels nutzt \autocite{the_gvisor_authors_security_2026}, wird gVisor als
User-Space-Kernel bezeichnet. Der Preis sind langsamere Systemaufrufe
\autocite[S.~3]{young_true_2019}.

Ein \emph{V8-Isolate} ist eine Instanz der JavaScript-Engine V8 mit eigenem Heap. Innerhalb eines
Isolates kann ein \emph{Context} voneinander unabhängigen JavaScript-Code in derselben Instanz
ausführen \autocite{the_v8_authors_getting_2026}. Die Grenze zwischen zwei Isolates verläuft damit
innerhalb eines Prozesses und wird allein von der Engine durchgesetzt. Sie ist eine sprachbasierte
Grenze, die Angriffe unterhalb der Softwareebene nicht abfängt, weshalb der Browser Chrome
Webseiten verschiedener Herkunft in getrennte Prozesse verlegt hat
\autocite[S.~1661--1662]{reis_site_2019}.

\emph{\texttt{deno\_core}} ist eine Rust-Bibliothek des Deno-Projekts, die den Umgang mit V8
kapselt. Ihre Runtime-Instanz hält ein Isolate und verbindet es mit einer Event Loop. Über
\emph{Ops} werden Rust-Funktionen so registriert, dass JavaScript sie aufrufen kann, und
\emph{Extensions} bündeln Ops mit dem JavaScript-Code, der sie dem Skript zugänglich macht
\autocite{iwanczuk_roll_2022}. Die Bibliothek bildet zugleich den Kern der Deno-Laufzeitumgebung
\autocite{deno_land_inc_how_2026}. \emph{Bindings} heißen in dieser Arbeit die Schnittstellen, über die ein
Skript die Außenwelt erreicht, etwa \texttt{fetch} oder ein Socket. Ein Binding besteht aus Ops
und dem JavaScript-Code, der sie dem Skript in der gewohnten Form anbietet.

\subsection{Bausteine der Ausführung}
\label{sec:grundlagen:bausteine}

\emph{Knative Serving} erweitert Kubernetes um Ressourcen für die Ausführung von Diensten, die auf
Anfragen reagieren. Ein Service verwaltet den Lebenszyklus eines Workloads, eine Revision ist der
Stand von Code und Konfiguration zu einem Zeitpunkt, und Revisionen werden nach eingehendem
Verkehr skaliert, bis hinunter auf null Pods \autocite{the_knative_authors_overview_2026}. Nach
welcher Wartezeit ohne Verkehr ein Pod entfernt wird, ist konfigurierbar
\autocite{the_knative_authors_configuring_2026}. Der Ausgangszustand der Fallstudie nutzt genau
dieses Verhalten, wie \autoref{sec:systemkontext} zeigt.

Ein \emph{Kaltstart} ist die Verzögerung, die eine Anfrage erfährt, wenn keine laufende Instanz sie
bedienen kann und die Plattform erst eine bereitstellen muss
\autocite[S.~1]{lin_mitigating_2019}. In der Fallstudie umfasst sie das Anlegen und Planen des
Pods, den Start der Laufzeitumgebung und das Laden des Codes. Eine Anfrage, die auf eine bereits
laufende Instanz trifft, gilt als \emph{warm}. Zwei Gegenmittel sind in der Literatur
beschrieben: ein Vorrat vorgehaltener Instanzen, der den Start vom Anfragepfad nimmt
\autocite[S.~1]{lin_mitigating_2019}, und Snapshots initialisierter Instanzen, aus denen ein Start
schneller hervorgeht als aus einem leeren Prozess \autocite[S.~559]{ustiugov_benchmarking_2021}.
Wie die Arbeit den Kaltstart misst, legt \autoref{sec:kriterienkatalog} fest.

\emph{NATS JetStream} ist die Persistenzschicht des Messaging-Systems NATS. Ein \emph{Stream}
speichert Nachrichten serverseitig für ein oder mehrere Subjects, ein \emph{Consumer} ist eine
serverseitige Sicht auf einen Stream, die festhält, wie weit ein Client ihn gelesen hat. Bestätigt
ein Client eine Nachricht nicht rechtzeitig, stellt der Server sie erneut zu
\autocite{the_nats_authors_jetstream_2026}. Die Plattform der Fallstudie betreibt einen solchen
Stream bereits für ihre Kommandos und Events, und das Ausführungsmodell nutzt ihn mit, wie
\autoref{sec:bausteine} beschreibt.

\subsection{Mandantenisolation und Zustandsfreiheit}
\label{sec:grundlagen:eigenschaften}

Ein \emph{Mandant} ist in der Fallstudie ein Kunde, dessen Nutzer sich eine Installation des
\ac{DMS} teilen. Das deckt sich mit der Definition von \textcite[S.~427]{krebs_architectural_2012},
nach der ein Mandant eine Gruppe von Nutzern mit derselben Sicht auf eine Anwendung ist und
Mandantenfähigkeit darin besteht, eine Anwendungsinstanz zwischen mehreren Mandanten zu teilen.
Mehrere Mandanten laufen auf derselben Plattform, und die Plattform muss sie
voneinander trennen, weil sie einander nicht vertrauen. \emph{Mandantenisolation} heißt in dieser
Arbeit, dass eine Ausführung Daten, Code und Verarbeitung anderer Mandanten weder einsehen noch
verändern noch stören kann. Welche Schutzziele daraus folgen und woran die Trennung heute
durchgesetzt wird, leitet \autoref{sec:systemkontext} aus dem Betriebskonzept her.

\emph{Zustandsfreiheit} bezieht sich nicht auf Mandanten, sondern auf aufeinanderfolgende
Ausführungen. Eine Ausführung ist zustandsfrei, wenn nichts von ihr einer späteren zugänglich ist,
weder im Speicher des Prozesses noch in globalen Variablen, registrierten Timern, offenen
Verbindungen oder im Dateisystem. Die Eigenschaft ist von der Mandantenisolation unabhängig, denn
sie kann auch innerhalb eines Mandanten verletzt werden, sobald eine Plattform Instanzen über
Ausführungen hinweg wiederverwendet, um Kaltstarts zu sparen. Warum diese Verletzung im
Agentenprofil zu einem Sicherheitsproblem wird, zeigt \autoref{sec:systemkontext}.
